% figures/tikz/board_options.tex --- where each area runs under the three options.
% % figures/tikz/board_options.tex --- where each area runs under the three options.
% % figures/tikz/board_options.tex --- where each area runs under the three options.
% % figures/tikz/board_options.tex --- where each area runs under the three options.
% \input{figures/tikz/board_options} inside a figure environment.
\begin{tikzpicture}[
    area/.style = {draw, minimum width=26mm, minimum height=6.5mm, font=\small},
    pi/.style   = {area, fill=blue!12},
    mcu/.style  = {area, fill=orange!22},
    hw/.style   = {area, fill=black!6, dashed},
    chip/.style = {font=\footnotesize\bfseries},
    head/.style = {font=\small\bfseries},
    link/.style = {{Latex[length=1.6mm]}-{Latex[length=1.6mm]}, thick},
  ]
  % Each column: five software areas stacked top to bottom, then the hardware.
  % #1 column x, #2..#6 style per area (pi or mcu)
  \newcommand{\optioncolumn}[7]{%
    \node[head] at (#1, 0.9) {#7};
    \node[#2] (c#1-1) at (#1, 0)    {Capture};
    \node[#3] (c#1-2) at (#1,-0.65) {Perception};
    \node[#4] (c#1-3) at (#1,-1.30) {Estimation};
    \node[#5] (c#1-4) at (#1,-1.95) {Navigation};
    \node[#6] (c#1-5) at (#1,-2.60) {Control};
    \node[hw] (c#1-6) at (#1,-3.45) {Actuation};
  }

  \optioncolumn{0}{mcu}{mcu}{mcu}{mcu}{mcu}{MCU-only}
  \optioncolumn{5}{pi}{pi}{pi}{pi}{mcu}{Pi + MCU}
  \optioncolumn{10}{pi}{pi}{pi}{pi}{pi}{Pi-only (chosen)}

  % The Pi + MCU boundary: the UART link between the two chips
  \draw[very thick, red!65!black, dashed] ([xshift=-3mm]c5-5.north west) -- ([xshift=3mm]c5-5.north east)
      node[right, font=\footnotesize, align=left, text=red!65!black] {UART\\link};

  % Legend
  \node[pi,  minimum width=6mm, label={[font=\footnotesize]right:Raspberry Pi}] at (1.6,-4.55) {};
  \node[mcu, minimum width=6mm, label={[font=\footnotesize]right:Microcontroller}] at (5.0,-4.55) {};
  \node[hw,  minimum width=6mm, label={[font=\footnotesize]right:Driver and motors}] at (8.6,-4.55) {};
\end{tikzpicture}
 inside a figure environment.
\begin{tikzpicture}[
    area/.style = {draw, minimum width=26mm, minimum height=6.5mm, font=\small},
    pi/.style   = {area, fill=blue!12},
    mcu/.style  = {area, fill=orange!22},
    hw/.style   = {area, fill=black!6, dashed},
    chip/.style = {font=\footnotesize\bfseries},
    head/.style = {font=\small\bfseries},
    link/.style = {{Latex[length=1.6mm]}-{Latex[length=1.6mm]}, thick},
  ]
  % Each column: five software areas stacked top to bottom, then the hardware.
  % #1 column x, #2..#6 style per area (pi or mcu)
  \newcommand{\optioncolumn}[7]{%
    \node[head] at (#1, 0.9) {#7};
    \node[#2] (c#1-1) at (#1, 0)    {Capture};
    \node[#3] (c#1-2) at (#1,-0.65) {Perception};
    \node[#4] (c#1-3) at (#1,-1.30) {Estimation};
    \node[#5] (c#1-4) at (#1,-1.95) {Navigation};
    \node[#6] (c#1-5) at (#1,-2.60) {Control};
    \node[hw] (c#1-6) at (#1,-3.45) {Actuation};
  }

  \optioncolumn{0}{mcu}{mcu}{mcu}{mcu}{mcu}{MCU-only}
  \optioncolumn{5}{pi}{pi}{pi}{pi}{mcu}{Pi + MCU}
  \optioncolumn{10}{pi}{pi}{pi}{pi}{pi}{Pi-only (chosen)}

  % The Pi + MCU boundary: the UART link between the two chips
  \draw[very thick, red!65!black, dashed] ([xshift=-3mm]c5-5.north west) -- ([xshift=3mm]c5-5.north east)
      node[right, font=\footnotesize, align=left, text=red!65!black] {UART\\link};

  % Legend
  \node[pi,  minimum width=6mm, label={[font=\footnotesize]right:Raspberry Pi}] at (1.6,-4.55) {};
  \node[mcu, minimum width=6mm, label={[font=\footnotesize]right:Microcontroller}] at (5.0,-4.55) {};
  \node[hw,  minimum width=6mm, label={[font=\footnotesize]right:Driver and motors}] at (8.6,-4.55) {};
\end{tikzpicture}
 inside a figure environment.
\begin{tikzpicture}[
    area/.style = {draw, minimum width=26mm, minimum height=6.5mm, font=\small},
    pi/.style   = {area, fill=blue!12},
    mcu/.style  = {area, fill=orange!22},
    hw/.style   = {area, fill=black!6, dashed},
    chip/.style = {font=\footnotesize\bfseries},
    head/.style = {font=\small\bfseries},
    link/.style = {{Latex[length=1.6mm]}-{Latex[length=1.6mm]}, thick},
  ]
  % Each column: five software areas stacked top to bottom, then the hardware.
  % #1 column x, #2..#6 style per area (pi or mcu)
  \newcommand{\optioncolumn}[7]{%
    \node[head] at (#1, 0.9) {#7};
    \node[#2] (c#1-1) at (#1, 0)    {Capture};
    \node[#3] (c#1-2) at (#1,-0.65) {Perception};
    \node[#4] (c#1-3) at (#1,-1.30) {Estimation};
    \node[#5] (c#1-4) at (#1,-1.95) {Navigation};
    \node[#6] (c#1-5) at (#1,-2.60) {Control};
    \node[hw] (c#1-6) at (#1,-3.45) {Actuation};
  }

  \optioncolumn{0}{mcu}{mcu}{mcu}{mcu}{mcu}{MCU-only}
  \optioncolumn{5}{pi}{pi}{pi}{pi}{mcu}{Pi + MCU}
  \optioncolumn{10}{pi}{pi}{pi}{pi}{pi}{Pi-only (chosen)}

  % The Pi + MCU boundary: the UART link between the two chips
  \draw[very thick, red!65!black, dashed] ([xshift=-3mm]c5-5.north west) -- ([xshift=3mm]c5-5.north east)
      node[right, font=\footnotesize, align=left, text=red!65!black] {UART\\link};

  % Legend
  \node[pi,  minimum width=6mm, label={[font=\footnotesize]right:Raspberry Pi}] at (1.6,-4.55) {};
  \node[mcu, minimum width=6mm, label={[font=\footnotesize]right:Microcontroller}] at (5.0,-4.55) {};
  \node[hw,  minimum width=6mm, label={[font=\footnotesize]right:Driver and motors}] at (8.6,-4.55) {};
\end{tikzpicture}
 inside a figure environment.
\begin{tikzpicture}[
    area/.style = {draw, minimum width=26mm, minimum height=6.5mm, font=\small},
    pi/.style   = {area, fill=blue!12},
    mcu/.style  = {area, fill=orange!22},
    hw/.style   = {area, fill=black!6, dashed},
    chip/.style = {font=\footnotesize\bfseries},
    head/.style = {font=\small\bfseries},
    link/.style = {{Latex[length=1.6mm]}-{Latex[length=1.6mm]}, thick},
  ]
  % Each column: five software areas stacked top to bottom, then the hardware.
  % #1 column x, #2..#6 style per area (pi or mcu)
  \newcommand{\optioncolumn}[7]{%
    \node[head] at (#1, 0.9) {#7};
    \node[#2] (c#1-1) at (#1, 0)    {Capture};
    \node[#3] (c#1-2) at (#1,-0.65) {Perception};
    \node[#4] (c#1-3) at (#1,-1.30) {Estimation};
    \node[#5] (c#1-4) at (#1,-1.95) {Navigation};
    \node[#6] (c#1-5) at (#1,-2.60) {Control};
    \node[hw] (c#1-6) at (#1,-3.45) {Actuation};
  }

  \optioncolumn{0}{mcu}{mcu}{mcu}{mcu}{mcu}{MCU-only}
  \optioncolumn{5}{pi}{pi}{pi}{pi}{mcu}{Pi + MCU}
  \optioncolumn{10}{pi}{pi}{pi}{pi}{pi}{Pi-only (chosen)}

  % The Pi + MCU boundary: the UART link between the two chips
  \draw[very thick, red!65!black, dashed] ([xshift=-3mm]c5-5.north west) -- ([xshift=3mm]c5-5.north east)
      node[right, font=\footnotesize, align=left, text=red!65!black] {UART\\link};

  % Legend
  \node[pi,  minimum width=6mm, label={[font=\footnotesize]right:Raspberry Pi}] at (1.6,-4.55) {};
  \node[mcu, minimum width=6mm, label={[font=\footnotesize]right:Microcontroller}] at (5.0,-4.55) {};
  \node[hw,  minimum width=6mm, label={[font=\footnotesize]right:Driver and motors}] at (8.6,-4.55) {};
\end{tikzpicture}
